% ===========================================================================
%  preamble.tex --- packages and macros for the birth--death--catastrophe
%  chapter.  Derived from the house preamble, CH2/preamble.tex; where the two
%  disagreed, Chapter 2 won.
%  When this chapter is dropped into the thesis, merge the contents of this
%  file into the thesis preamble; duplicates may simply be deleted.
% ===========================================================================

\usepackage[T1]{fontenc}
\usepackage[utf8]{inputenc}
\usepackage{lmodern}
\usepackage{microtype}
\usepackage{amsmath,amssymb,amsthm,mathtools}
\usepackage{graphicx}
\usepackage{float}
\usepackage{placeins}
\usepackage{booktabs}
\usepackage[hypcap=false]{caption}
\usepackage{tikz}
\usetikzlibrary{arrows.meta,positioning,calc}
\usepackage{pgfplots}
\pgfplotsset{compat=1.18}
\usepackage{enumitem}
\usepackage{xcolor}
\usepackage{array}
\usepackage{hyperref}
\usepackage[nameinlink,noabbrev]{cleveref}

\graphicspath{{figures/}}

% --- Equation cross-references as bare parenthesised numbers -----------------
\crefformat{equation}{(#2#1#3)}
\Crefformat{equation}{(#2#1#3)}
\crefrangeformat{equation}{(#3#1#4)--(#5#2#6)}
\crefmultiformat{equation}{(#2#1#3)}{ and (#2#1#3)}{, (#2#1#3)}{ and (#2#1#3)}

% --- Notation ---------------------------------------------------------------
% Scalars are set in italic; bold is reserved for vector-valued states, which
% appear only in the multi-type work of later chapters.
\newcommand{\pgf}{probability generating function}
\newcommand{\rv}{random variable}
\newcommand{\pr}[1]{\mathbb{P}\!\left(#1\right)}
\newcommand{\ex}[1]{\mathbb{E}\!\left(#1\right)}
\newcommand{\var}[1]{\operatorname{Var}\!\left(#1\right)}
\newcommand{\cova}[1]{\operatorname{Cov}\!\left(#1\right)}
\newcommand{\dd}{\mathrm{d}}
\newcommand{\ee}{\mathrm{e}}
\newcommand{\ddt}[1]{\frac{\dd #1}{\dd t}}
\newcommand{\dda}[2]{\frac{\dd #1}{\dd #2}}
\newcommand{\pddt}[1]{\frac{\partial #1}{\partial t}}
\newcommand{\pdd}[2]{\frac{\partial #1}{\partial #2}}
\newcommand{\lrb}[1]{\left(#1\right)}
\newcommand{\lrbs}[1]{\left[#1\right]}
\newcommand{\ind}{\mathbf 1}

% --- Chapter-specific notation ----------------------------------------------
% \Ifix is the non-fixation probability; its optional argument carries the
% founder count, as in \Ifix[,k].  \Kb is the burst size; note that K(t) is
% the second moment of the load, a different object.
\newcommand{\Ifix}[1][]{I_{\mathrm{fix}#1}}
\newcommand{\Kb}{\mathcal{K}}
\newcommand{\yp}{\textit{Y.~pestis}}
\newcommand{\ft}{\textit{F.~tularensis}}

% --- Cross-chapter references -----------------------------------------------
% Backward references to the mathematical introduction.  Redefining these is
% all that thesis integration requires.
\newcommand{\ChM}{the mathematical introduction}
\newcommand{\ChMcap}{The mathematical introduction}

% Readable standalone fallbacks for labels defined in the mathematical
% introduction.  Copied from CH2/sections/02_markov_chains.tex, which uses the
% same device for labels defined in Section 1.3.  In the assembled thesis the
% real cleveref target is used; compiled alone, the second argument prints.
\providecommand{\mextref}[2]{%
  \ifcsname r@#1\endcsname\cref{#1}\else#2\fi}

% --- Forward references to later modelling chapters -------------------------
% Every forward reference is routed through \fwd, whose first argument is a
% stable key and whose second is the prose actually typeset.  Keys in use
% here: singlecell (the PGF/QSD/burst-law chapter), population (the renewal
% population chapter).
\newcommand{\fwd}[2]{#2}

% --- Theorem environments ---------------------------------------------------
\theoremstyle{plain}
\newtheorem{theorem}{Theorem}[chapter]
\newtheorem{proposition}[theorem]{Proposition}
\newtheorem{lemma}[theorem]{Lemma}
\newtheorem{corollary}[theorem]{Corollary}
\theoremstyle{definition}
\newtheorem{definition}[theorem]{Definition}
\theoremstyle{remark}
\newtheorem{remark}[theorem]{Remark}

\allowdisplaybreaks
